\subsection{CH2021 smoothings and limits}
\label{subsec:generalized-reeb-orbits-polytopes-ch2021}

Chaidez--Hutchings study a different finite model for four-dimensional
polytopes, designed to compare Reeb dynamics on a polytope with Reeb dynamics
on smoothings of that polytope \cite{CH2021}. Their main hypotheses and
conclusions are not the same as the Haim--Kislev formula above. The
Haim--Kislev formula applies to convex polytopes in all even dimensions and
computes the EHZ capacity by a finite quadratic program. The
Chaidez--Hutchings construction is four-dimensional and uses a
\emph{symplectic polytope} hypothesis, which in particular excludes
Lagrangian two-faces so that the relevant combinatorial Reeb direction is
one-dimensional.

For such a four-dimensional polytope, a Chaidez--Hutchings combinatorial Reeb
orbit is a cyclic sequence of oriented line segments in the boundary. Each
segment lies in a face and points in the corresponding one-dimensional Reeb
cone. They distinguish three types. Type~1 orbits avoid the one-skeleton,
Type~2 orbits meet the one-skeleton only at finitely many segment endpoints,
and Type~3 orbits contain a bad one-face. The symplectic flow graph records
the Type~1 case: vertices are two-faces, edges are Reeb passages through
three-faces, and periodic graph orbits correspond to Type~1 combinatorial Reeb
orbits with the same action \cite[Proposition~2.2]{CH2021}.

The limit results in \cite{CH2021} explain why this combinatorial model is
useful for capacity computations on four-dimensional symplectic polytopes.
First, a nondegenerate Type~1 combinatorial Reeb orbit gives nearby Reeb orbits
on sufficiently small smoothings; their actions converge, while the rotation
number and Conley--Zehnder index agree with the combinatorial values for all
sufficiently small smoothing parameters \cite[Theorem~1.8]{CH2021}. Conversely,
given Reeb orbits on \(\partial X_{\varepsilon_i}\) with
\(\varepsilon_i\to0\) and uniformly bounded rotation number, a subsequence
converges to a Type~1 or Type~2 combinatorial Reeb orbit, again with action
convergence \cite[Theorem~1.9]{CH2021}. Combining this with the fact that an
EHZ-capacity minimizer can be chosen with rotation number at most \(2\), they
obtain a capacity formula for four-dimensional symplectic polytopes: minimize
over Type~2 combinatorial Reeb orbits and over Type~1 combinatorial Reeb orbits
with combinatorial rotation number at most \(2\), subject in both cases to
their finite-segment bound \cite[Corollary~1.13]{CH2021}.

This thesis uses these statements as background for the later flow-graph
section, not as a replacement for the Haim--Kislev quadratic program. The
current QP checkpoint uses the Haim--Kislev formula as the capacity interface.
The flow-graph algorithm based on CH2021, including its implemented scope and
its unsupported cases, is treated separately in
Section~\ref{sec:flow-graph-algorithm-ch2021}.
